\chapter*{Úvod}
\addcontentsline{toc}{chapter}{Úvod}

\section*{Vědecký obor}

Fyzika studuje obecné vlastnosti látek a~polí. Fyzika je \emph{vědecký obor}, a~proto všechna tvrzení v~ní musejí být přímo nebo nepřímo, alespoň v~principu, vyvratitelná experimentem (falzifikovatelná).

Konstrukce fyzikálních teorií se obecně řídí tzv.~\emph{Occamovou břitvou}: modely a tvrzení by měly být co nejjednodušší, avšak v~souladu se všemi dostupnými pozorováními. Není-li nějaká část teorie nezbytná pro vysvětlení pozorovaných jevů, z~teorie se odstraní.

Fyzika a chemie mají široký překryv. Atomy a molekuly se řídí zákony kvantové mechaniky, i~když odvození zákonů chemie přímo ze základních zákonů fyziky představuje silně netriviální problém mnoha částic.

\subsection*{Mechanika}

Hrubé, nejednoznačné a nepřesné, zato však běžně používané dělení fyzikálních oborů je například:
\begin{description}
    \item[mechanika] pohyb těles a jejich vzájemné silové působení,
    \item[termodynamika a statistická fyzika] jevy způsobené chaotickým pohybem velkého množství částic,
    \item[elektřina a magnetismus] elektromagnetické pole a jeho interakce s~hmotou,
    \item[jaderná a částicová fyzika] jevy v~atomovém jádře a struktura mikrosvěta.
\end{description}

Mechanika je obor fyziky, který se zabývá pohybem těles a jejich vzájemným působením. Samotným popisem pohybu bez ohledu na jeho příčiny se zabývá \emph{kinematika}. Příčiny pohybu, tedy proč pohyb vzniká nebo se mění v~čase vlivem sil, jsou předmětem \emph{dynamiky}.

V~prvním přiblížení se budeme věnovat pohybu \emph{hmotných bodů}, tj.~idealizaci reálných objektů, které lze plně popsat pouze jejich hmotností, polohou a rychlostí. Hmotné body nemají žádnou vnitřní strukturu ani vlastní rozměry. \emph{Klasická (newtonovská)} mechanika popisuje pohyb těles, která jsou dostatečně těžká ve srovnání s~atomy, ale ne natolik hmotná, aby bylo nutné uvažovat zakřivení prostoročasu, a která se na makroskopických škálách pohybují rychlostmi podstatně menšími, než je rychlost světla.

Newtonovská mechanika selhává, když se rychlosti těles začnou blížit rychlosti světla (čímž se dostáváme do \emph{speciální teorie relativity}) nebo pokud jsou gravitační pole příliš silná (což je doména \emph{obecné teorie relativity}).

Naopak při malých hmotnostech a na mikroskopických vzdálenostech se dostáváme do oblasti \emph{kvantové mechaniky}, kde klasické pojmy přesné polohy a trajektorie částice ztrácejí svůj původní smysl.